\appendix

\chapter{Codebase Implementations and Algorithmic Blueprints}\label{cha:appendix_code}

Code snippets implementing ideas described in Chapters~\ref{cha:chapter5}, \ref{cha:chapter6}, and \ref{cha:chapter7}, excluding standard boilerplate code and logging routines.

\section{Dynamic SDF World Parsing and Coordinate Geometry}\label{app:code_sdf_parser}

The code implementing the 2D rotation matrix transformations and CAD origin offset compensation is given below in Listing~\ref{lst:sdf_transform_core}. It reproduces the calculations described in Section~5.3.2.

\begin{lstlisting}[language=Python, caption={Coordinate transformation matrix and geometric center offset resolution.}, label={lst:sdf_transform_core}]
def apply_transformation_matrix(
    raw_x: float, 
    raw_y: float, 
    yaw: float, 
    local_dx: float, 
    local_dy: float
) -> Tuple[float, float]:
    """
    Applies a 2D rotation matrix to map local offsets into global world coordinates.
    Formula:
        gx = raw_x + local_dx * cos(yaw) - local_dy * sin(yaw)
        gy = raw_y + local_dx * sin(yaw) + local_dy * cos(yaw)
    """
    cos_yaw = math.cos(yaw)
    sin_yaw = math.sin(yaw)
    global_x = raw_x + (local_dx * cos_yaw) - (local_dy * sin_yaw)
    global_y = raw_y + (local_dx * sin_yaw) + (local_dy * cos_yaw)
    return round(global_x, 3), round(global_y, 3)

def get_item_geometric_center(
    item_type: str, 
    raw_x: float, 
    raw_y: float, 
    yaw: float
) -> Tuple[float, float]:
    """Calculates the true geometric center of an item by applying its local offset."""
    if item_type in ("shelf", "shelf_big"):
        # Shelves have -0.5m offset along local X in Gazebo models
        local_dx = -0.5
        local_dy = 0.0
    else:
        local_dx = 0.0
        local_dy = 0.0
        
    return apply_transformation_matrix(raw_x, raw_y, yaw, local_dx, local_dy)
\end{lstlisting}

Also in Section~5.3.2, Listing~\ref{lst:sdf_shelving_core} shows the implementation of the dynamic waypoints generation algorithm for both small-scale retail shelf and steel column bypass.

\begin{lstlisting}[language=Python, caption={Parametric section layout generation and steel-column structural bypass.}, label={lst:sdf_shelving_core}]
def build_small_shelf_sides(name: str, cx: float, cy: float, yaw: float) -> Dict[str, Any]:
    """Generates sides and section waypoints for a standard small shelf."""
    sides = {}
    labels = ["A", "B", "C", "D"]
    half_width = SMALL_SHELF_SPECS["width"] / 2.0
    step = SMALL_SHELF_SPECS["section_step"]
    offset_dist = (SMALL_SHELF_SPECS["depth"] / 2.0) + APPROACH_DISTANCE_M # 1.5m
    
    side_1_pts, side_2_pts = [], []
    for i in range(len(labels)):
        local_dx = -half_width + (step / 2.0) + (i * step)
        side_1_pts.append(apply_transformation_matrix(cx, cy, yaw, local_dx, offset_dist))
        side_2_pts.append(apply_transformation_matrix(cx, cy, yaw, local_dx, -offset_dist))
        
    side_1_y_avg = sum(pt[1] for pt in side_1_pts) / len(side_1_pts)
    side_2_y_avg = sum(pt[1] for pt in side_2_pts) / len(side_2_pts)
    
    yplus_raw, yminus_raw = (side_1_pts, side_2_pts) if side_1_y_avg > side_2_y_avg else (side_2_pts, side_1_pts)
        
    for side_name, raw_pts in [("yplus", yplus_raw), ("yminus", yminus_raw)]:
        raw_pts.sort(key=lambda pt: pt[0]) # Ascending X order
        section_points_list = [(pt[0], pt[1], labels[idx]) for idx, pt in enumerate(raw_pts)]
        sides[side_name] = {
            "side": side_name,
            "key": f"{name}_{side_name}",
            "approach_x": section_points_list[0][0],
            "approach_y": section_points_list[0][1],
            "sections": SMALL_SHELF_SPECS["sections"],
            "section_points_template": section_points_list,
            "section_labels_template": labels,
        }
    return sides

# --- Steel Column Structural Bypass Snippet in parse_sdf() ---
structural_shifts = [("shelf_big_0", "xplus"), ("shelf_big_4", "xminus")]
for shelf_name, side_name in structural_shifts:
    if shelf_name in items and side_name in items[shelf_name]["sides"]:
        side_spec = items[shelf_name]["sides"][side_name]
        shifted_points = []
        for gx, gy, label in side_spec["section_points_template"]:
            if label == "C":
                gy = -13.50  # Safely offset Y coordinate to clear steel column at Y = -15.00m
            shifted_points.append((gx, gy, label))
        side_spec["section_points_template"] = shifted_points
\end{lstlisting}

\section{Peer-to-Peer Spatial Mutex and Claim Manager}\label{app:code_claim_manager}

The code of the decentralized spatial mutual exclusion mechanism, including the TTL lease management is implemented in the class \texttt{ClaimManager} from \texttt{waypoint\_sender.py}. Listing~\ref{lst:claim_manager_core} reproduces the relevant snippets for claim processing, LCFS lease overwrite procedure, and blacklist propagation described in Sections~\ref{sec:ch5_claims} and~\ref{sec:ch6_blacklist}.\footnote{Corresponding to Section~5.2.4 and Section~6.5.2 in Chapter~5 and Chapter~6.}

\begin{lstlisting}[language=Python, caption={Peer-to-peer claim manager and lease maintenance engine.}, label={lst:claim_manager_core}]
class ClaimManager:
    def __init__(self, node: Node):
        self.node = node
        self.active_claims: Dict[str, Dict[str, Any]] = {}
        self.local_active_claim_key: Optional[str] = None
        
        qos_claims = QoSProfile(depth=100, reliability=ReliabilityPolicy.RELIABLE)
        self.claims_pub = self.node.create_publisher(String, "/side_claims", qos_claims)

    def process_claim_message(self, side_key: str, robot_id: str, action: str, expire_at_ns: int):
        """Updates internal claims registry based on global claims messages."""
        if action == "release":
            self.active_claims.pop(side_key, None)
            if robot_id == self.node.robot_name:
                self.local_active_claim_key = None
        elif action in ("claim", "renew"):
            self.active_claims[side_key] = {
                "robot_id": robot_id,
                "expire_at_ns": expire_at_ns
            }
            if robot_id == self.node.robot_name:
                self.local_active_claim_key = side_key
        elif action == "blacklist":
            if robot_id != self.node.robot_name:
                self.node.permanently_blocked_side_keys.add(side_key)
                self.node.deferred_side_keys.discard(side_key)
                self.active_claims.pop(side_key, None)

    def publish_local_claim(self, action_string: str):
        """Publishes own claims or renewals to the global topic."""
        if not self.node.active_target_dict:
            return
            
        now_ns = int(self.node.get_clock().now().nanoseconds)
        expire_at_ns = now_ns + (6 * 1_000_000_000) # 6.0 second lease TTL
        
        side_key = self.node.active_target_dict['side_key']
        payload = {
            "robot_id": self.node.robot_name,
            "side_key": side_key,
            "action": action_string,
            "timestamp_ns": now_ns,
            "expire_at_ns": expire_at_ns
        }
        
        msg = String()
        msg.data = json.dumps(payload)
        self.claims_pub.publish(msg)

    def release_local_claim(self):
        """Releases local spatial claim and notifies the peer fleet."""
        if not self.local_active_claim_key:
            return
            
        now_ns = int(self.node.get_clock().now().nanoseconds)
        payload = {
            "robot_id": self.node.robot_name,
            "side_key": self.local_active_claim_key,
            "action": "release",
            "timestamp_ns": now_ns,
            "expire_at_ns": now_ns
        }
        
        msg = String()
        msg.data = json.dumps(payload)
        self.claims_pub.publish(msg)
        self.active_claims.pop(self.local_active_claim_key, None)
        self.local_active_claim_key = None
\end{lstlisting}

\section{Progress-Based Kinematic Stall Watchdog and Escape Engine}\label{app:code_recovery}

The progress-based motion stall detection algorithm and the deterministic open-loop physical escape are implemented in \texttt{recovery.py}. Listing~\ref{lst:recovery_manager_core} below shows the $10\mathrm{Hz}$ tick function which evaluates the progress delta ($\delta_{\text{progress}} = 0.03\mathrm{m}$ over $10.0\mathrm{s}$), initiates an open-loop linear backup ($v_x = -0.12\mathrm{m/s}$, $1.67\mathrm{s}$), and the direction-alternating rotational escape ($\omega_z = \pm 0.35\text{rad/s}$, $0.9\mathrm{s}$). These parameters are determined by the calculations described in Section~6.3.

\begin{lstlisting}[language=Python, caption={Progress-based stall detection watchdog and physical escape engine.}, label={lst:recovery_manager_core}]
class RecoveryManager:
    def __init__(self, node: Node, cmd_vel_publisher, nav2_cancel_cb, nav2_retrigger_cb, defer_task_cb):
        self.node = node
        self.cmd_vel_pub = cmd_vel_publisher
        self.cancel_nav2_goal = nav2_cancel_cb
        self.retrigger_nav2_goal = nav2_retrigger_cb
        self.defer_active_task = defer_task_cb
        
        self.state = RecoveryState.IDLE
        self.retry_count = 0
        self.last_movement_timestamp_ns = 0
        self.maneuver_end_timestamp_ns = 0
        
        self.PROGRESS_EPSILON_M = 0.03  # Must move 3cm to reset stall clock
        self.BACKWARD_SPEED = -0.12
        self.BACKWARD_DURATION_S = 0.2 / 0.12  # 1.667s duration to cover 0.20m
        self.ESCAPE_ROTATION_SPEED = 0.35
        self.ESCAPE_ROTATION_DURATION_S = 0.9

    def tick(self, current_x: float, current_y: float):
        """Master 10Hz clock tick of the recovery manager."""
        if self.state == RecoveryState.IDLE:
            return
            
        now_ns = self.node.get_clock().now().nanoseconds
        
        if self.state == RecoveryState.MONITORING:
            elapsed_seconds = (now_ns - self.last_movement_timestamp_ns) / 1e9
            threshold = STALL_CANDIDATE_TIMEOUT_SEC if self.is_heading_to_candidate else STALL_DIRECT_TIMEOUT_SEC
            if elapsed_seconds > threshold:
                self._handle_stall_detection()
                
        elif self.state == RecoveryState.BACKING_UP:
            if now_ns >= self.maneuver_end_timestamp_ns:
                self._stop_cmd_vel_output()
                self._start_escape_rotation()
            else:
                self._publish_backward_drive()
                
        elif self.state == RecoveryState.ROTATING_ESCAPE:
            if now_ns >= self.maneuver_end_timestamp_ns:
                self._stop_cmd_vel_output()
                self._trigger_nav2_replan()
            else:
                self._publish_escape_rotation()

    def _start_physical_escape(self):
        if self.retry_count >= 5:
            self.state = RecoveryState.EXHAUSTED
            self.defer_active_task()
            return

        self.retry_count += 1
        self.state = RecoveryState.BACKING_UP
        now_ns = self.node.get_clock().now().nanoseconds
        self.maneuver_end_timestamp_ns = now_ns + int(self.BACKWARD_DURATION_S * 1e9)

    def _publish_escape_rotation(self):
        cmd = Twist()
        # Alternate rotation direction based on attempt number to clear costmap edges
        rotation_sign = 1.0 if (self.retry_count % 2 == 1) else -1.0
        cmd.angular.z = rotation_sign * self.ESCAPE_ROTATION_SPEED
        self.cmd_vel_pub.publish(cmd)
\end{lstlisting}

\section{Onboard Edge Perception, Analytical Redundancy, and Linear XAI Back-Projection}\label{app:code_xai}

The onboard visual perception module and the linear XAI attribution model are implemented in \texttt{waypoint\_sender.py}. Listing~\ref{lst:xai_backproject_core} shows the reproduction of the closed-form vision back-projection algorithm ($\mathbf{w}_{\text{raw}} = \mathbf{w}_{\text{pca}} \mathbf{V}$), unscaling procedure, $4 \times 4$ spatial grid interpolation, and alpha-blending of the colormap ($\alpha_{\text{blend}} = 0.35$) from Section~5.6.

\begin{lstlisting}[language=Python, caption={Closed-form linear SVM weight back-projection and XAI colormap generation.}, label={lst:xai_backproject_core}]
def _generate_xai_overlay(self, perturbed_img):
    """Generates mathematically honest Explainable AI (XAI) overlays based on SVM coefficients."""
    if len(perturbed_img.shape) == 2:
        perturbed_color = cv2.cvtColor(perturbed_img, cv2.COLOR_GRAY2BGR)
    else:
        perturbed_color = perturbed_img.copy()

    if not self.scaler or not self.pca or not self.svm_active or not hasattr(self.svm_model, "coef_"):
        return perturbed_color

    try:
        # 1. Project Linear SVM decision boundary weights back to 1924-dim space (W_raw = W_pca * V^T)
        pca_weights = self.svm_model.coef_[0]
        raw_weights = np.dot(pca_weights, self.pca.components_)
        
        # 2. Rescale using StandardScaler's array to reverse scaling influence
        raw_weights_scaled = raw_weights / (self.scaler.scale_ + 1e-8)
        
        # 3. Extract the first 160 dimensions (Uniform Spatial LBP features)
        lbp_weights_magnitude = np.abs(raw_weights_scaled[:160])
        
        # 4. Sum 10 bins across each of the 16 blocks to calculate spatial grid importance
        block_importances = lbp_weights_magnitude.reshape(16, 10).sum(axis=1)
        grid_importance = block_importances.reshape(4, 4)
        
        # 5. Normalize and upscale 4x4 grid to 128x128 via bilinear interpolation
        grid_min, grid_max = grid_importance.min(), grid_importance.max()
        grid_norm = ((grid_importance - grid_min) / (grid_max - grid_min + 1e-8) * 255).astype(np.uint8)
        heatmap = cv2.resize(grid_norm, (128, 128), interpolation=cv2.INTER_LINEAR)
        
        # 6. Apply JET Colormap for visual overlay
        heatmap_color = cv2.applyColorMap(heatmap, cv2.COLORMAP_JET)
        
        # 7. Blend 65% original camera view with 35% XAI decision boundary heatmap
        overlay = cv2.addWeighted(perturbed_color, 0.65, heatmap_color, 0.35, 0)
        return overlay

    except Exception as e:
        self.get_logger().error(f"[{self.robot_name}] Failed to generate XAI heatmap: {e}")
        return perturbed_color
\end{lstlisting}

\section{Deterministic Launch Orchestration and Readiness Gates}\label{app:code_launch}

The deterministic robot launch procedure requires two types of modifications to the code. First, the blocking readiness gates are defined in \texttt{gate.py}, while the launch process itself is orchestrated in \texttt{multi\_robot\_controlled.launch.py} using \texttt{OnProcessExit} event handlers. The code snippet in Listing~\ref{lst:gate_primitives_core} below reproduces the relevant blocks implementing the temporal constraints from Section~7.2.2.

\begin{lstlisting}[language=Python, caption={Deterministic readiness gate primitives for launch orchestration.}, label={lst:gate_primitives_core}]
def wait_for_clock(poll_sec: float = 1.0) -> None:
    """Block until at least one message arrives on /clock."""
    print('[gate:clock] waiting for /clock...', flush=True)
    while True:
        result = _run(['ros2', 'topic', 'echo', '--once', '/clock', '--field', 'clock.sec'], timeout=3.0)
        if result and result.returncode == 0 and result.stdout.strip():
            print('[gate:clock] ready - /clock is publishing', flush=True)
            return
        time.sleep(poll_sec)

def wait_for_lifecycle(nodes: list, poll_sec: float = 1.0) -> None:
    """Block until every lifecycle node in `nodes` reports state active."""
    assert nodes, '[gate:lifecycle] nodes list must not be empty'
    while True:
        not_active = []
        for node in nodes:
            result = _run(['ros2', 'lifecycle', 'get', node], timeout=4.0)
            if not result or result.returncode != 0:
                not_active.append(node)
                continue
            out = result.stdout.strip().lower()
            if not (out.startswith('active') or 'active [3]' in out):
                not_active.append(node)
        if not not_active:
            print('[gate:lifecycle] ready - all nodes active', flush=True)
            return
        time.sleep(poll_sec)
\end{lstlisting}

The next step in deterministic launch orchestration is to define event-driven \texttt{OnProcessExit} handlers for each robot in the system to ensure sequential launch without race-conditions between robots. This is implemented in the \texttt{multi\_robot\_controlled.launch.py} module, and the code snippet is given below in Listing~\ref{lst:launch_chains_core}. It reproduces the code blocks described in Section~7.2.3.

\begin{lstlisting}[language=Python, caption={Sequential per-robot lifecycle daisy-chaining via OnProcessExit event handlers.}, label={lst:launch_chains_core}]
def _robot_chain(robot, world_name, base_urdf, bridge_yaml, nav2_yaml, map_yaml, default_bt_xml, start_trigger_gate):
    name = robot['name']
    x, y = robot['x'], robot['y']

    spawn = ExecuteProcess(
        cmd=[sys.executable, _SPAWN_PY, '--world-name', world_name, '--robot-name', name,
             '--sdf-file', robot['sdf'], '--x', x, '--y', y, '--z', '0.01'],
        name=f'spawn_{name}', output='screen'
    )

    topics_gate = _gate('topics', topics=[f'/{name}/scan', f'/{name}/odom_filtered', f'/{name}/joint_states', '/tf', '/tf_static'], label=f'gate_topics_{name}')
    lifecycle_gate = _gate('lifecycle', nodes=[f'/{name}/map_server', f'/{name}/amcl', f'/{name}/planner_server', f'/{name}/controller_server', f'/{name}/behavior_server', f'/{name}/bt_navigator'], label=f'gate_lifecycle_{name}')

    spawn_trigger = RegisterEventHandler(OnProcessExit(target_action=start_trigger_gate, on_exit=[spawn]))

    actions_list = [
        spawn_trigger,
        RegisterEventHandler(OnProcessExit(target_action=spawn, on_exit=[bridge, rsp, ekf_node, topics_gate])),
        RegisterEventHandler(OnProcessExit(target_action=topics_gate, on_exit=[_nav2_group(name, nav2_yaml, map_yaml, default_bt_xml, x, y), lifecycle_gate])),
        RegisterEventHandler(OnProcessExit(target_action=lifecycle_gate, on_exit=_mission_nodes(name))),
    ]
    return actions_list, lifecycle_gate
\end{lstlisting}

\section{System Configuration Blueprints and Narrow-Aisle Mutex Allocations}\label{app:code_config}

The code snippet below in Listing~\ref{lst:config_blueprints_core} reproduces the implementation of the physical geometry constants, static allocation tables, and narrow-aisle spatial mutexes from \texttt{config.py}.

\begin{lstlisting}[language=Python, caption={Global physical constants and narrow aisle mutual exclusion mappings.}, label={lst:config_blueprints_core}]
APPROACH_DISTANCE_M = 1.2
NAVIGATION_FINAL_GOAL_TOLERANCE_M = 0.35

SMALL_SHELF_SPECS = {
    "width": 3.6, "depth": 0.6, "sections": 4, "section_step": 0.9,
    "offset_local_x": -0.5, "offset_local_y": 0.0
}

BIG_SHELF_SPECS = {
    "width": 2.1, "length": 18.0, "sections": 4, "section_step": 4.5,
    "offset_local_x": -0.5, "offset_local_y": 0.0
}

# --- Narrow Aisle Mutual Exclusion Groups (Corridor Mutex) ---
SHARED_CORRIDORS = {
    "shelf_big_0_xplus": "corridor_west",
    "shelf_big_4_xminus": "corridor_west",
    "shelf_big_4_xplus": "corridor_mid",
    "shelf_big_3_xminus": "corridor_mid",
    "shelf_big_3_xplus": "corridor_east",
    "shelf_big_2_xminus": "corridor_east",
    "shelf_2_yminus": "corridor_small_west_top",
    "shelf_1_yplus": "corridor_small_west_top",
    "shelf_9_yplus": "corridor_small_corner_top",
    "shelf_10_yminus": "corridor_small_corner_top"
}
\end{lstlisting}